\chapter{Stridor}

\section*{Definition}
Stridor is a high-pitched, harsh, monophonic breath sound resulting from turbulent airflow through a narrowed or obstructed upper airway (pharynx, larynx, trachea, or mainstem bronchi). It is a musical, predominantly inspiratory sound, though biphasic or expiratory stridor may occur depending on the level and fixity of the obstruction.

\section*{Pathophysiology}
Turbulent airflow through a partially obstructed upper airway generates the characteristic sound of stridor. Bernoulli's principle states that flow velocity increases at a constriction, reducing lateral wall pressure and potentially collapsing the segment (especially in extrathoracic airways during inspiration). The timing of stridor helps localize the lesion: inspiratory stridor suggests supraglottic or glottic pathology; biphasic stridor indicates a fixed tracheal or subglottic obstruction; expiratory stridor suggests intrathoracic tracheal obstruction.

\section*{Etiology}
\begin{itemize}
  \item Infectious causes: Croup (laryngotracheobronchitis -- most common in children), Epiglottitis (Haemophilus influenzae type b, Streptococcus spp.), Bacterial tracheitis, Retropharyngeal abscess, Peritonsillar abscess, Diphtheria, Laryngeal papillomatosis
  \item Congenital causes (pediatric): Laryngomalacia (most common congenital laryngeal anomaly), Vocal cord paralysis, Subglottic stenosis, Vascular rings, Laryngeal web, Tracheoesophageal fistula
  \item Inflammatory/immune: Anaphylaxis, Angioedema (allergic or hereditary), Laryngospasm, Granulomatosis with polyangiitis, Sarcoidosis, Rheumatoid arthritis (cricoarytenoid involvement)
  \item Neoplastic causes: Laryngeal or tracheal tumors (squamous cell carcinoma), Thyroid tumors with tracheal invasion, Mediastinal masses, Esophageal carcinoma
  \item Traumatic/iatrogenic: Post-intubation subglottic stenosis, Foreign body aspiration, Neck trauma, Vocal cord hematoma, Post-thyroidectomy recurrent laryngeal nerve injury
\end{itemize}

\section*{Indications for Evaluation}
Seek care if:
\begin{itemize}
  \item Any stridor is present, particularly in children (airway may deteriorate rapidly)
  \item Acute onset of stridor with drooling, dysphonia, or tripod positioning (suspect epiglottitis)
  \item Stridor with respiratory distress: retractions, nasal flaring, cyanosis, altered mental status
  \item Post-traumatic or post-extubation onset of stridor
  \item Suspected foreign body aspiration in children or adults
\end{itemize}

\section*{Related Symptoms}
\begin{itemize}
  \item Dysphonia
  \item Hoarseness
  \item Cough (barking/croupy in croup)
  \item Dysphagia
  \item Drooling
  \item Respiratory distress
  \item Cyanosis
  \item Tachycardia
  \item Hypoxia
\end{itemize}
\textbf{Note:} Stridor is a medical emergency requiring immediate assessment of airway patency; bedside flexible laryngoscopy can localize the obstruction but should not delay securing the airway when distress is evident.

\section*{Self-Care / Home Management}
\begin{itemize}
  \item Acute stridor is a medical emergency -- do not attempt home management. Call 911 or present to the nearest emergency department.
  \item For mild croup at home, cool mist or steam inhalation may provide temporary relief while seeking medical assessment.
  \item Keep the child calm and upright to minimize airflow turbulence and crying which worsens obstruction.
  \item Avoid placing anything in the mouth, including attempting to visualize the throat, with suspected epiglottitis.
  \item Do not delay transport for medical evaluation while waiting for symptoms to improve.
\end{itemize}

\section*{Diagnostic Approach}
\begin{itemize}
  \item Rapid ABC assessment (airway, breathing, circulation) with pulse oximetry.
  \item Flexible nasolaryngoscopy at bedside or in controlled setting to visualize supraglottic and glottic structures.
  \item Neck X-ray (anteroposterior and lateral) for epiglottitis, croup (steeple sign), or retropharyngeal abscess.
  \item CT neck/chest for deep space infections, neoplasm, or external compression.
  \item Direct laryngoscopy and bronchoscopy under anesthesia for definitive diagnosis and therapeutic intervention.
\end{itemize}

\section*{Treatment / Management Overview}
\begin{itemize}
  \item Airway management: Supplemental oxygen, heliox (helium-oxygen mixture to reduce turbulent flow), or endotracheal intubation/tracheostomy for impending failure.
  \item Croup: Nebulized epinephrine (racemic or L-epinephrine) and systemic corticosteroids (dexamethasone).
  \item Epiglottitis: Broad-spectrum antibiotics (ceftriaxone) and airway protection; intubation in operating room with ENT standby.
  \item Anaphylaxis/angioedema: Epinephrine IM, antihistamines, corticosteroids, and airway monitoring.
  \item Foreign body: Rigid bronchoscopy for removal.
  \item Subglottic stenosis: Endoscopic dilation, laser resection, or open surgical reconstruction.
\end{itemize}

\clearpage
